\documentclass{article}

\usepackage[english]{babel}
\usepackage[letterpaper,top=2cm,bottom=2cm,left=3cm,right=3cm,marginparwidth=1.75cm]{geometry}
\usepackage{amsmath,amsfonts}
\usepackage{graphicx}
\usepackage{booktabs}
\usepackage{array}
\usepackage[left]{lineno}

% 行号样式，可选
\renewcommand\linenumberfont{\normalfont\tiny\sffamily\color{gray}}
\setlength\linenumbersep{6pt}
\usepackage{tabularx}
\usepackage{longtable}
\usepackage{pdflscape}
\usepackage{float}
\usepackage{xcolor}
\usepackage[section]{placeins}
\usepackage[colorlinks=true,allcolors=blue]{hyperref}

\renewcommand{\thefigure}{S\arabic{figure}}
\renewcommand{\thetable}{S\arabic{table}}
\renewcommand{\theHfigure}{S\arabic{figure}}
\renewcommand{\theHtable}{S\arabic{table}}

\title{Supplementary Material for Hippocampal Controllability across Clinically Defined Alzheimer's Disease-Related Groups: Associations with Structural Network Topology and Group Discrimination}
\author{}
\date{}
\modulolinenumbers[1]
\linenumbers

\begin{document}
\maketitle

\section{Supplementary Methods}

Throughout the Supplementary Material, BH-FDR-adjusted \(p\) values are denoted by \(p_{\mathrm{FDR}}\), whereas unadjusted and permutation \(p\) values are denoted by \(p\) and \(p_{\mathrm{perm}}\), respectively.

\subsection{Cortical and Structural-Connectivity Visualizations}
For descriptive cortical visualization, ROI-level average and modal controllability values were averaged within each clinical group and mapped to the Schaefer2018 cortical surface using a common color scale (Supplementary Fig.~S1). Group-averaged structural connectivity matrices were also used to visualize hippocampal connections and connections between Tian atlas-defined non-cortical structures and the seven Schaefer networks. Only relatively high-weight connections were retained for display; this visualization threshold was not used in any inferential analysis. Mean connectivity between each atlas-defined structure group and cortical network was summarized in heatmaps (Supplementary Fig.~S2).

\subsection{Classification Analyses}
The four binary tasks were CN versus AD, SMC versus MCI, MCI versus AD, and CN versus SMC. Six low-dimensional feature sets were examined: global, left, and right hippocampal controllability; bilateral FreeSurfer hippocampal volume; hippocampal strength and betweenness; and their six-feature combination. Controllability predictors were within-subject ranks, volume and topology were unranked, and the combined model contained both scales. The 24 task--feature combinations were exploratory and formed one Holm-corrected family. All models used inverse-frequency-balanced \(L_2\)-regularized logistic regression and the same covariates: age, sex, and education; ICV was not a classification covariate. Ten repetitions of five-fold stratified outer cross-validation generated held-out probabilities, while four-fold stratified inner cross-validation selected \(C\in\{0.01,0.1,1,10\}\). Covariate residualization and standardization were fitted within each inner-training partition during hyperparameter selection and refitted on the corresponding outer-training data before held-out prediction. Each participant's 10 held-out probabilities were averaged arithmetically before calculation of one subject-level AUC.

Classifier significance used 1,000 complete-pipeline label permutations for every task--feature combination. Labels were permuted before cross-validation; every permutation reran the full repeated nested procedure, including foldwise residualization, standardization, inner selection of \(C\), model refitting, and outer prediction. Permutation \(p\) values used the plus-one correction. AUC confidence intervals used 5,000 group-stratified bootstrap resamples of the fixed subject-level mean out-of-fold probabilities. These intervals quantify subject-sampling uncertainty conditional on the fitted predictions and do not fully capture training-set variability; the mean, standard deviation, and range of the 10 repeat-specific AUCs are reported in Table~II. Direct candidate-versus-volume comparisons reran both models in the same complete-case sample using identical splits and the same random seed as Table~II. Paired AUC differences used 5,000 group-stratified bootstrap resamples and subject-level score-swap tests; the 20 comparisons formed a separate Holm-corrected family (Supplementary Table~\ref{tab:paired_volume_classification}). The high-dimensional connectivity analysis remains exploratory; its previous fixed-score permutation values are removed and only descriptive AUC estimates are retained.

\subsection{Primary Rank Endpoints and Unranked-Value Sensitivity Analyses}
The primary endpoints were mean within-subject ranks: nodal average and modal controllability were ranked separately from 1 to 450 within each participant and then averaged across the bilateral hippocampal parcels. Adjusted group means, 95\% confidence intervals, and partial \(\eta^2\) were obtained from the full group-plus-demographics models using 5,000 group-stratified bootstrap resamples. Six pairwise contrasts per endpoint were tested with two-sided Freedman--Lane permutations and controlled jointly as one 12-test BH-FDR family.

Unranked bilateral hippocampal values were analyzed after adjustment for age, sex, education, and \(\log(1+W_{\mathrm{total}})\), where \(W_{\mathrm{total}}=\sum_{i<j}A_{ij}\) and \(A\) was the final inverse-node-volume-scaled connectivity matrix before spectral stabilization. A further model added connectome density. Unranked omnibus and density-adjusted omnibus tests formed separate two-outcome BH-FDR families, and the 12 unranked pairwise contrasts formed another family. Direct average--modal associations used partial Spearman correlations: rank-based measures were adjusted for age, sex, and education, whereas unranked measures were additionally adjusted for matrix scale.

\subsection{Cohort Construction and Diagnostic Sensitivity Analysis}
The cohort was reconstructed from the ADNI imaging-metadata export used for the analysis, the final subject manifest, and the executable audit code. A session was a unique participant-date combination. Same-date duplicate diffusion series were resolved using the lower numeric image identifier, and the earliest eligible date was retained when multiple dates remained. The previous 240 participants were retained; systematic application of the same eligibility criteria identified 9 additional SMC and 32 additional AD participants. The archived materials supported a reconstructed query inventory and partial cohort-selection audit. They did not preserve the historical intermediate exclusion reasons; no retrospective reason was assigned without an archived record (Supplementary Table~\ref{tab:cohort_flow_audit}).

Clinical-group assignment in the primary analysis followed the source-cohort labels for all 281 participants. Diagnostic records temporally matched to the selected diffusion-MRI sessions were examined in sensitivity analyses, which did not alter the interpretation of the primary group effects.

\subsection{Cognitive Associations and Education Moderation}
Partial Spearman associations between the rank-based bilateral hippocampal average- and modal-controllability measures and MMSE, ADAS-Cog 11, ADAS-Cog 13,
CDR memory, and MoCA were estimated after residualizing cognition and
controllability for age, sex, and education. For direct comparison with
hippocampal volume, the analysis was repeated in the same FreeSurfer-complete
participants after additional adjustment for ICV. BH-FDR correction was
applied within each defined association family. Education was retained
as a primary covariate because it differed among groups, but its potential role
as a cognitive-reserve proxy was examined in separate
diagnosis-by-education interaction models. The three interaction terms were
tested jointly using 10,000 Freedman--Lane permutations.

\subsection{Vascular and White-Matter-Hyperintensity Analyses}
Modified Hachinski Ischemic Scale and UCD white-matter-hyperintensity (WMH) measurements were matched to the selected diffusion MRI session within 365 days. WMH burden was calculated as \(100\times\mathrm{TOTAL\_WMH}/\mathrm{CEREBRUM\_TCV}\) and transformed using \(\log(1+x)\). Group differences were assessed using Kruskal--Wallis tests. Partial Spearman associations with hippocampal controllability, topology, WMH burden, and cognition controlled for clinical group, age, sex, and education; BH-FDR correction was applied separately within each vascular predictor and outcome family.

\subsection{MCI Subtype Audit}
The ADNI \texttt{DXSUM} record nearest to the selected diffusion MRI session was used to audit MCI subtype. The primary window was 180 days, with 365-day and any-record audits as sensitivity checks. Memory-feature MCI was defined from \texttt{DXMDES}; a stricter aMCI definition additionally required \texttt{DXMPTR5}=Yes, indicating objective memory impairment for age and education. The primary clinical-group models were repeated after restricting the MCI group to these definitions. These classifications were used only for sensitivity analysis and do not establish subsequent conversion to AD dementia.

\subsection{Amyloid PET Sensitivity Analysis}
The selected diffusion MRI session was matched to the nearest visually quality-control-passed ADNI UC Berkeley 6-mm amyloid PET record within 365 days. Binary amyloid status and continuous Centiloid burden were analyzed in separate complete-case models because their availability differed; valid negative Centiloid values were retained. Models included diagnosis, age, sex, education, and either amyloid variable. Diagnosis effects adjusted for amyloid and amyloid effects adjusted for diagnosis were evaluated using 10,000 Freedman--Lane permutations. The analysis was repeated with a 730-day window. Quality-control-passed flortaucipir PET coverage was audited, but the four-group tau analysis was not interpreted because only one SMC participant had an eligible scan within 365 days.

\subsection{Hippocampal Long-Axis Analysis}
The Tian S3 atlas contains medial-head, lateral-head, body, and tail parcels in each hemisphere. Homologous parcels were averaged to obtain four bilateral long-axis measures from within-subject ranked controllability values. The primary clinical-group model included age, sex, and education as
covariates and used 10,000 Freedman--Lane permutations. BH-FDR correction
was applied across the four subdivisions separately for average and
modal controllability. These macroscopic connectivity-based parcels are not equivalent to cytoarchitectonic hippocampal subfields.

\subsection{Shared and Incremental Information}
The source was the ADNI UCSF Cross-Sectional FreeSurfer 7.x export dated 11 July 2026. ICV, left hippocampal volume, and right hippocampal volume were \texttt{ST10CV}, \texttt{ST29SV}, and \texttt{ST88SV}. Eligible records were 3.0~T, had \texttt{EXAMDATE} exactly equal to the selected diffusion-MRI date, and contained all three measurements; same-day duplicates were resolved using the lower numeric \texttt{IMAGEUID}. All 270 included records used FreeSurfer 7.2.0. \texttt{STATUS} was \texttt{partial}, and \texttt{OVERALLQC} and \texttt{HIPPOQC} were not populated; volume-containing analyses were therefore exploratory.

Shared-information models predicted unranked average and modal controllability from age, sex, education, ICV, left and right hippocampal volumes entered as separate predictors, strength, and betweenness. Incremental cognitive models used the same base predictors and jointly added rank-based average and modal controllability. Reduced and full models used identical complete cases. The five two-degree-of-freedom partial \(F\)-tests formed one BH-FDR family. Percentile 95\% confidence intervals for changes in \(R^2\) used 5,000 group-stratified bootstrap resamples, with reduced and full models refitted in every resample. Because \(\Delta R^2\) is constrained to be nonnegative in nested ordinary least-squares models, these intervals were interpreted descriptively; inferential conclusions were based on the partial \(F\)-tests, with BH-FDR correction applied across the five outcomes. Clinical-group membership was analyzed using multinomial logistic regression, with a 5,000-resample percentile interval for the change in McFadden \(R^2\).

\subsection{Multiplicity Families}
Multiplicity was controlled within defined analysis families rather than across every analysis in the study. The family definitions, number of tests, and correction procedures are summarized in Supplementary Table~\ref{tab:multiplicity_families}. Unless otherwise specified, adjusted values use BH-FDR; the exploratory classification families use Holm familywise-error correction.

\subsection{Global Gray-Matter Sensitivity Analysis}
Selected diffusion MRI sessions were matched to same-day ADNI UCSF Cross-Sectional FreeSurfer 7.x records. Non-hippocampal gray matter was calculated by subtracting bilateral hippocampal volume from total gray matter. Partial Spearman associations were estimated after residualizing for ICV, age, sex, and education, then repeated with clinical group included. Duplicate same-day records were resolved using the lower image identifier. Because the FreeSurfer records lacked final visual-QC ratings, these analyses were considered exploratory.

\subsection{Diffusion-Protocol and Registration QC Audit}
\begin{sloppypar}
The archived processing environment was MRtrix3 3.0.6 and FSL 6.0.7.18, including FLIRT. Five-tissue-type images were generated with \texttt{5ttgen fsl}. Response functions were estimated with \texttt{dwi2response dhollander}, and the same \texttt{dwi2fod msmt\_csd} command, with white-matter, gray-matter, and CSF outputs, was applied to all datasets. This included 266 datasets with one nonzero shell and 15 with three nonzero shells; the use of the same multi-tissue command in the single-nonzero-shell datasets is therefore retained as a processing limitation. The final DWI inputs contained nine total-volume categories (22, 31, 36, 42, 46, 54, 55, 68, and 127), b0-count categories of 1, 4, 5, 6, 7, 8, and 13, and nonzero diffusion-weighted volume-count categories of 16, 30, 32, 36, 41, 48, 60, and 114. Observed shell sets were $b=0/1000$, $b=0/1300$, and $b=0/500/1000/2000$~s/mm$^{2}$. Fifteen 127-volume inputs were PA acquisitions, each containing 13 b0 and 114 nonzero diffusion-weighted volumes. Protocol-specific counts are given in Supplementary Table~\ref{tab:protocol_qc_summary}.
\end{sloppypar}

For retrospective registration QC, the mean b0 image was registered to the first T1-derived 5TT volume using six-degree-of-freedom FLIRT with mutual-information cost and trilinear interpolation. The transform was imported with \texttt{transformconvert}; its inverse was applied to the T1-derived brain mask in the mean-b0 grid with linear interpolation. The combined Schaefer--Tian atlas was transferred from MNI to native T1 space using six-degree-of-freedom FLIRT with mutual-information cost and nearest-neighbour label interpolation. Direct overlay inspection confirmed that the resulting native-T1 node image and tractogram shared the same physical coordinate frame. The node image was therefore supplied directly to \texttt{tck2connectome}, without an additional voxel-grid resampling step before streamline-to-node assignment. For retrospective QC only, the native-T1 atlas was transformed into the mean-b0 grid with nearest-neighbour interpolation to calculate coverage and generate overlays.

Let $M_{\mathrm{DWI}}$ and $M_{\mathrm{T1}}$ denote the DWI brain mask and transformed T1 brain mask, and let $H$ denote the transformed bilateral hippocampal-label mask. Brain-mask overlap was quantified as $2|M_{\mathrm{DWI}}\cap M_{\mathrm{T1}}|/(|M_{\mathrm{DWI}}|+|M_{\mathrm{T1}}|)$, and hippocampal coverage as $|H\cap M_{\mathrm{DWI}}|/|H|$. Dice below 0.70 or hippocampal coverage below 0.90 triggered targeted review. Registration montages displayed transformed T1 brain-mask contours over the mean b0 image, and hippocampal montages displayed transformed hippocampal-label contours over the same image. These measures identify gross mask-overlap or out-of-brain placement problems but do not establish anatomical localization accuracy. The audit was completed quantitatively for all 281 participants. The five coverage-trigger cases comprised one CN, two MCI, and two AD participants. Manual inspection confirmed acceptable T1--DWI registration and hippocampal-label placement in all five cases, without gross misregistration or erroneous hippocampal localization. Their exclusion was evaluated in a two-endpoint sensitivity family. Within the final 281-participant analytic cohort, there were no additional registration, tractography/connectome-availability, or QC exclusions. Candidate-level failure counts preceding formation of the final analytic cohort were not available in the archived processing records.

\subsection{Motion and Multicenter Sensitivity Analyses}
Mean and maximum absolute and relative RMS displacement were derived from each participant's FSL \texttt{eddy\_movement\_rms} output. Group differences were tested using Kruskal--Wallis tests, and the primary controllability models were repeated with mean absolute and relative RMS displacement as additional covariates. Multicenter sensitivity models sequentially adjusted for scanner manufacturer and ADNI phase; site; and, in the strict model, scanner model, phase, site, acquisition acceleration, NIfTI-derived diffusion-volume-count category, slice thickness, and both motion measures. After reconciliation against the final NIfTI inputs, the strict model contained 12 scanner-model, three phase, 43 site, two acceleration, nine volume-count, and four slice-thickness levels; its nominal 76-column full design had rank 71 and 210 residual degrees of freedom. All controllability models used 10,000 Freedman--Lane permutations. ComBat was not applied because diagnosis, scanner, phase, and site were not fully crossed.

Acquisition overlap was evaluated in three complementary analyses. First, generalized overlap weights were estimated from scanner model, phase, site, acceleration, NIfTI-derived volume-count category, slice thickness, and mean absolute and relative RMS displacement. Balance was assessed using the maximum absolute pairwise standardized mean difference and group-specific effective sample sizes. Second, a connected exact-cell common-support analysis retained exact acquisition cells containing at least two clinical groups, with the additional requirement that every pair of clinical groups shared at least one cell. The model adjusted for age, sex, education, both motion measures, and acquisition-cell indicators; Freedman--Lane residual permutations were restricted within cell. Third, we quantified the feasibility of the still stricter requirement that every retained cell contain all four clinical groups. Because overlap weighting retained material imbalance and no exact cell contained all four groups after protocol reconciliation, these diagnostics and the connected common-support result were interpreted jointly rather than as proof that acquisition confounding had been removed.

\subsection{SMC Sample-Size Sensitivity Analyses}
Balanced subsampling was repeated 1,000 times with 27 participants per group. The partial \(F\), \(\Delta R^2\), significance indicator, and adjusted SMC direction were recorded in each repetition. Each SMC participant was also removed in turn, and the primary 10,000-permutation model was rerun. Because resamples overlapped, these analyses were interpreted as stability checks rather than independent replication.

\section{Supplementary Results}

\subsection{Cortical and Connectivity Visualizations}
The descriptive cortical maps were broadly similar across groups, and none of the seven aggregated cortical networks survived BH-FDR correction. Visual variation in Supplementary Fig.~S1 therefore should not be interpreted as parcel-level statistical evidence. Supplementary Fig.~S2 shows widespread hippocampal and Tian atlas-defined non-cortical structural connectivity with cortical networks; these thresholded displays were descriptive only.

\subsection{Rank and Unranked Hippocampal Results}
Both matrix-scale-adjusted unranked omnibus effects were retained, but the average-controllability pairwise pattern was scale dependent. The SMC--CN and SMC--MCI differences survived correction only for within-subject ranks, whereas the AD--MCI difference survived correction only in the matrix-scale-adjusted unranked analysis (Supplementary Tables~\ref{tab:primary_pairwise_contrasts} and \ref{tab:raw_pairwise}).

\subsection{MCI Subtype Results}
Within 180 days of the selected diffusion-MRI session, 80 of the 88 MCI participants had sufficient records for subtype classification. Of these, 76 had memory-only features, two had mixed memory and non-memory features, and two had non-memory-only features. Seventy-seven met the stricter memory-feature-plus-objective-impairment definition. Clinical-group effects remained significant after restriction to either amnestic definition (all \(p_{\mathrm{FDR}}\leq0.0012\); Supplementary Table~\ref{tab:tableS1}).

\subsection{Hippocampal Long-Axis Results}
Clinical-group effects survived correction in the bilateral medial head, lateral head, body, and tail for both controllability measures, indicating macroscopic long-axis heterogeneity without resolving histological subfields (Supplementary Table~\ref{tab:tableS2}).

\subsection{Education Moderation Results}
Diagnosis-by-education interactions were not significant for average or modal controllability (both \(p_{\mathrm{FDR}}=0.748\); Supplementary Table~\ref{tab:tableS3}).

\subsection{Cognitive Association Results}
Demographic-adjusted controllability--cognition associations were significant after BH-FDR correction (absolute \(\rho=0.134\text{--}0.256\)). In the matched ICV-adjusted comparison, hippocampal volume showed larger associations with cognition (absolute \(\rho=0.225\)--0.453) than controllability (absolute \(\rho=0.124\)--0.253; Supplementary Table~\ref{tab:tableS4}).

\subsection{Vascular and WMH Results}
Modified Hachinski scores did not differ among groups and were not associated with the analyzed imaging or cognitive outcomes after correction. Normalized WMH burden differed among groups and, after diagnosis and demographic adjustment, was associated with lower average controllability (\(\rho=-0.142\), \(p_{\mathrm{FDR}}=0.044\)), higher modal controllability (\(\rho=0.137\), \(p_{\mathrm{FDR}}=0.044\)), and lower MoCA (\(\rho=-0.247\), \(p_{\mathrm{FDR}}=0.006\); Supplementary Table~\ref{tab:tableS5}).

\subsection{Shared and Incremental Information Results}
In the 270-participant exact-date FreeSurfer sample, the harmonized structural models explained 54.6\% and 55.0\% of raw average- and modal-controllability variance. Adding rank-based average and modal controllability produced an FDR-corrected increment only for CDR memory among the five cognitive outcomes (\(\Delta R^2=0.0249\), 95\% CI 0.0047--0.0638, \(p_{\mathrm{FDR}}=0.0374\)). The MMSE increment was nominal before correction but did not survive BH-FDR correction (\(\Delta R^2=0.0149\), 95\% CI 0.0000--0.0530, \(p=0.0418\), \(p_{\mathrm{FDR}}=0.0974\)). The multinomial model retained \(\Delta R^2_{\mathrm{McFadden}}=0.0219\) (95\% CI 0.0103--0.0606, \(p=0.0189\); Supplementary Table~\ref{tab:tableS6}).

\subsection{Motion Sensitivity Results}
None of the four RMS motion summaries differed significantly among groups (all \(p_{\mathrm{FDR}}\geq0.233\)), and group effects remained significant after motion adjustment (Supplementary Table~\ref{tab:tableS7}).

\subsection{Multicenter Sensitivity Results}
Clinical-group effects remained significant in the strict scanner/site/protocol/motion model after reconciliation to the final NIfTI volume counts (average, \(p_{\mathrm{FDR}}=0.0093\); modal, \(p_{\mathrm{FDR}}=0.0093\); Supplementary Table~\ref{tab:tableS8}). The full design had rank 71 and 210 residual degrees of freedom. Generalized overlap weighting left a maximum absolute pairwise standardized mean difference of 1.042 and an overall effective sample size of 45.2 (group-specific values, 34.8 CN, 11.5 SMC, 6.7 MCI, and 20.9 AD), indicating inadequate residual balance; weighted estimates were therefore retained only as diagnostics. The connected exact-cell common-support analysis retained 218 participants in 46 cells (CN 93, SMC 24, MCI 66, and AD 35). Group effects remained significant for average rank (\(F_{3,164}=5.663\), \(p_{\mathrm{FDR}}=0.0025\), partial \(\eta^2=0.094\)) and modal rank (\(F_{3,164}=5.112\), \(p_{\mathrm{FDR}}=0.0025\), partial \(\eta^2=0.086\)). No exact acquisition cell contained all four groups. Direct SMC--MCI overlap comprised four cells containing five SMC and six MCI participants; the analysis therefore does not establish that the SMC--MCI classifier is free of acquisition confounding (Supplementary Table~\ref{tab:acquisition_overlap}). The temporally matched diagnostic-record sensitivity analysis did not alter the interpretation of either primary endpoint.

\subsection{Registration and Label-Coverage QC Results}
The retrospective audit was completed for all 281 participants. Five cases crossed the hippocampal-coverage review threshold (one CN, no SMC, two MCI, and two AD). Manual inspection confirmed acceptable registration and hippocampal-label placement in all five cases. Excluding these cases retained the average-rank group effect (\(F_{3,269}=5.554\), \(p_{\mathrm{FDR}}=0.0017\), partial \(\eta^2=0.058\)) and modal-rank group effect (\(F_{3,269}=5.264\), \(p_{\mathrm{FDR}}=0.0017\), partial \(\eta^2=0.055\); Supplementary Table~\ref{tab:protocol_qc_summary}).

\subsection{SMC Stability Results}
Balanced subsampling retained the group effect in 99.9\% of average- and 99.7\% of modal-controllability repetitions, and both effects remained significant in all leave-one-SMC-out analyses (Supplementary Table~\ref{tab:tableS9}).

\subsection{Global Gray-Matter Results}
In 270 FreeSurfer-matched participants, hippocampal volume covaried with total and non-hippocampal gray matter after ICV and demographic adjustment and after additional diagnosis adjustment (Supplementary Table~\ref{tab:tableS10}).

\subsection{Amyloid PET Results}
Within 365 days, binary amyloid status was available for 230 participants and Centiloids for 226. Diagnosis effects remained significant after adjustment for either measure, whereas neither amyloid measure was independently associated with controllability after diagnosis adjustment. The 730-day analysis gave the same conclusion (Supplementary Table~\ref{tab:tableS11}).

\subsection{Classification Results}
Complete low-dimensional results are reported in Table~II. Eleven of 24 complete-pipeline permutation tests remained significant after Holm correction. Bilateral volume had the highest AUC point estimates for CN versus AD (0.820) and MCI versus AD (0.825). For SMC versus MCI, the combined feature set had the highest estimate (0.729), while left controllability had the highest estimate among individual feature sets (0.701). In the common 112-person sample, none of the SMC-versus-MCI candidate-versus-volume differences survived Holm correction (Supplementary Table~\ref{tab:paired_volume_classification}). The classification results therefore describe task-specific point-estimate ordering without evidence that controllability is statistically superior to volume. The high-dimensional connectivity AUCs are retained as descriptive estimates only (Supplementary Table~\ref{tab:tableS12}).

\section{Supplementary Tables}

\begin{table}[!htbp]
\centering
\caption{MCI subtype audit and aMCI-restricted sensitivity analyses.}
\label{tab:tableS1}
\footnotesize
\setlength{\tabcolsep}{3.5pt}
\renewcommand{\arraystretch}{1.10}
\begin{tabularx}{\textwidth}{@{}>{\raggedright\arraybackslash}Xcccccc@{}}
\toprule
\textbf{Matching window or model} &
\textbf{Usable subtype record} &
\textbf{Memory only} &
\textbf{Mixed} &
\textbf{Non-memory only} &
\textbf{Strict aMCI} &
\textbf{Unclassified}\\
\midrule
\multicolumn{7}{l}{\textbf{A. MCI subtype audit}}\\
Within 180 days & 80 & 76 & 2 & 2 & 77 & 8\\
Within 365 days & 84 & 80 & 2 & 2 & 81 & 4\\
Any recorded MCI visit & 88 & 83 & 2 & 3 & 84 & 0\\
\bottomrule
\end{tabularx}

\vspace{3mm}
\begin{tabularx}{\textwidth}{@{}>{\raggedright\arraybackslash}Xlccccc@{}}
\toprule
\textbf{MCI inclusion model} &
\textbf{Measure} &
\textbf{\(n_{\mathrm{MCI}}\)} &
\textbf{Total \(n\)} &
\textbf{\(F\)} &
\textbf{\(p_{\mathrm{perm}}\)} &
\textbf{\(p_{\mathrm{FDR}}\)}\\
\midrule
\multicolumn{7}{l}{\textbf{B. Clinical-group sensitivity analyses}}\\
All clinically defined MCI
 & Average & 88 & 281 & 6.495 & 0.0004 & 0.0007\\
 & Modal   & 88 & 281 & 6.041 & 0.0007 & 0.0007\\
\addlinespace
Memory-feature aMCI within 180 days
 & Average & 78 & 271 & 6.485 & 0.0011 & 0.0012\\
 & Modal   & 78 & 271 & 6.061 & 0.0012 & 0.0012\\
\addlinespace
Memory feature + objective impairment
 & Average & 77 & 270 & 6.584 & 0.0004 & 0.0006\\
 & Modal   & 77 & 270 & 6.181 & 0.0006 & 0.0006\\
\bottomrule
\end{tabularx}

\vspace{1mm}
\footnotesize{\textit{Note:} The first numeric column reports availability of the records required for subtype classification within 180 days of the selected diffusion-MRI session. Memory only denotes an amnestic-only \texttt{DXMDES} record; mixed denotes both memory and non-memory features. Strict aMCI required a memory-feature \texttt{DXMDES} record and \texttt{DXMPTR5}=Yes, indicating objective memory impairment for age and education. Panel B reports Freedman--Lane permutation \(F\)-tests with 10,000 permutations adjusted for age, sex, and education. BH-FDR correction was applied across average and modal controllability within each inclusion model. These subtype analyses do not establish subsequent conversion to AD dementia.}
\end{table}

\begin{table}[!htbp]
\centering
\caption{Clinical-group effects across Tian S3 hippocampal long-axis subdivisions.}
\label{tab:tableS2}
\footnotesize
\setlength{\tabcolsep}{5pt}

\begin{tabular}{@{}llcccc@{}}
\toprule
\textbf{Metric} &
\textbf{Bilateral subdivision} &
\textbf{$n$} &
\textbf{$F(3,274)$} &
\textbf{$p_{\mathrm{perm}}$} &
\textbf{$p_{\mathrm{FDR}}$} \\
\midrule
Average & Medial head  & 281 & 6.712 & 0.0007 & 0.0024 \\
Average & Lateral head & 281 & 5.693 & 0.0012 & 0.0024 \\
Average & Body         & 281 & 3.825 & 0.0111 & 0.0148 \\
Average & Tail         & 281 & 3.387 & 0.0183 & 0.0183 \\
\midrule
Modal  & Medial head  & 281 & 6.566 & 0.0002 & 0.0008 \\
Modal  & Lateral head & 281 & 5.009 & 0.0025 & 0.0050 \\
Modal  & Body         & 281 & 3.816 & 0.0099 & 0.0132 \\
Modal  & Tail         & 281 & 2.735 & 0.0455 & 0.0455 \\
\bottomrule
\end{tabular}

\vspace{1mm}

\footnotesize{\textit{Note:} Omnibus clinical-group effects were tested after
adjustment for age, sex, and education using 10,000 Freedman--Lane
permutations. BH-FDR correction was applied across the four subdivisions
separately within average and modal controllability. Tian S3 long-axis
parcels are macroscopic connectivity-based subdivisions and are not
equivalent to histological hippocampal subfields.}

\end{table}

\begin{table}[!htbp]
\centering
\caption{Exploratory diagnosis-by-education interaction analyses for hippocampal controllability.}
\label{tab:tableS3}
\footnotesize
\setlength{\tabcolsep}{4pt}
\begin{tabular}{@{}lccccc@{}}
\toprule
\textbf{Outcome} & \textbf{$n$} & \textbf{$F(3,271)$} & \textbf{$p_{\mathrm{perm}}$} & \textbf{$p_{\mathrm{FDR}}$} & \textbf{$\Delta R^2$}\\
\midrule
Average controllability & 281 & 0.415 & 0.737 & 0.748 & 0.004 \\
Modal controllability & 281 & 0.403 & 0.748 & 0.748 & 0.004 \\
\bottomrule
\end{tabular}

\vspace{1mm}
\footnotesize{\textit{Note:} Models included age, sex, education, diagnosis, and the three diagnosis-by-education interaction terms, with CN as the reference group. The reported \(F\)-test compares models with and without the interaction block using 10,000 Freedman--Lane permutations. Education and age were standardized. BH-FDR correction was applied across the two controllability outcomes.}

\end{table}

\begin{table}[!htbp]
\centering
\caption{Partial Spearman correlations of hippocampal controllability and volume with cognitive performance.}
\label{tab:tableS4}
\footnotesize
\setlength{\tabcolsep}{3pt}

\textbf{A. Primary controllability--cognition associations}

\vspace{1mm}
\begin{tabular}{@{}lccc@{}}
\toprule
\textbf{Cognitive measure} & \textbf{$n$} & \textbf{Average controllability} & \textbf{Modal controllability}\\
 & & \textbf{$\rho$ ($p_{\mathrm{FDR}}$)} & \textbf{$\rho$ ($p_{\mathrm{FDR}}$)}\\
\midrule
MMSE & 279 & \(0.146\;(0.021)\) & \(-0.134\;(0.025)\) \\
ADAS-Cog 11 & 264 & \(-0.189\;(0.005)\) & \(0.180\;(0.006)\) \\
ADAS-Cog 13 & 264 & \(-0.194\;(0.005)\) & \(0.180\;(0.006)\) \\
CDR memory & 280 & \(-0.256\;(<0.001)\) & \(0.243\;(<0.001)\) \\
MoCA & 169 & \(0.183\;(0.021)\) & \(-0.176\;(0.025)\) \\
\bottomrule
\end{tabular}

\vspace{3mm}
\textbf{B. Matched ICV-adjusted comparison with hippocampal volume}

\vspace{1mm}
\begin{tabular}{@{}lcccc@{}}
\toprule
\textbf{Cognitive measure} & \textbf{$n$} & \textbf{Average controllability} & \textbf{Modal controllability} & \textbf{Hippocampal volume}\\
 & & \textbf{$\rho$ ($p_{\mathrm{FDR}}$)} & \textbf{$\rho$ ($p_{\mathrm{FDR}}$)} & \textbf{$\rho$ ($p_{\mathrm{FDR}}$)}\\
\midrule
MMSE & 268 & \(0.139\;(0.024)\) & \(-0.124\;(0.043)\) & \(0.367\;(<0.001)\) \\
ADAS-Cog 11 & 254 & \(-0.196\;(0.003)\) & \(0.183\;(0.005)\) & \(-0.437\;(<0.001)\) \\
ADAS-Cog 13 & 254 & \(-0.204\;(0.002)\) & \(0.187\;(0.005)\) & \(-0.453\;(<0.001)\) \\
CDR memory & 269 & \(-0.253\;(<0.001)\) & \(0.236\;(<0.001)\) & \(-0.393\;(<0.001)\) \\
MoCA & 166 & \(0.195\;(0.015)\) & \(-0.188\;(0.018)\) & \(0.225\;(0.005)\) \\
\bottomrule
\end{tabular}

\vspace{1mm}
\footnotesize{\textit{Note:} Panel A residualized the cognitive and controllability variables
for age, sex, and education. BH-FDR correction was applied across the 10
controllability--cognition associations in Panel A. Panel B used the same
participants for all three imaging measures and additionally adjusted for ICV;
BH-FDR correction was applied across its 15 comparisons. Higher MMSE and
MoCA scores indicate better cognition, whereas higher ADAS-Cog and CDR memory
scores indicate greater impairment. Hippocampal volume is the sum of the left
and right FreeSurfer 7.x volumes.}

\end{table}

\begin{table}[!htbp]
\centering
\caption{Associations of Modified Hachinski score and
white-matter hyperintensity burden with hippocampal imaging features
and cognition.}
\label{tab:tableS5}
\scriptsize
\setlength{\tabcolsep}{3.5pt}
\renewcommand{\arraystretch}{1.08}

\begin{tabularx}{\textwidth}{@{}>{\raggedright\arraybackslash}X
                                  >{\raggedright\arraybackslash}p{3.5cm}
                                  cccc@{}}
\toprule
\textbf{Analysis or predictor} & \textbf{Outcome or coverage} &
\textbf{\(n\)} & \textbf{Statistic} & \textbf{\(p\)} &
\textbf{\(p_{\mathrm{FDR}}\)}\\
\midrule

\multicolumn{6}{l}{\textbf{A. Data coverage}}\\
CN      & Hachinski / WMH / both & 110 & 88 / 110 / 88  & -- & --\\
SMC     & Hachinski / WMH / both & 27  & 21 / 27 / 21   & -- & --\\
MCI     & Hachinski / WMH / both & 88  & 67 / 87 / 67   & -- & --\\
AD      & Hachinski / WMH / both & 56  & 52 / 54 / 51   & -- & --\\
Overall & Hachinski / WMH / both & 281 & 228 / 278 / 227 & -- & --\\
\addlinespace

\multicolumn{6}{l}{\textbf{B. Clinical-group differences}}\\
Modified Hachinski score & Four-group comparison
 & 228 & \(H=4.329\) & 0.228 & 0.228\\
Normalized WMH burden & Four-group comparison
 & 278 & \(H=27.490\) & \(<0.0001\) & \(<0.0001\)\\
\addlinespace

\multicolumn{6}{l}{\textbf{C. Diagnosis- and demographic-adjusted partial Spearman associations}}\\
Modified Hachinski score & Average controllability
 & 228 & \(\rho=0.071\) & 0.286 & 0.692\\
Modified Hachinski score & Modal controllability
 & 228 & \(\rho=-0.063\) & 0.346 & 0.692\\
Modified Hachinski score & Hippocampal strength
 & 228 & \(\rho=0.039\) & 0.560 & 0.746\\
Modified Hachinski score & Hippocampal betweenness
 & 228 & \(\rho=0.009\) & 0.893 & 0.893\\
Modified Hachinski score & MMSE
 & 227 & \(\rho=-0.084\) & 0.206 & 0.846\\
Modified Hachinski score & ADAS-Cog 11
 & 212 & \(\rho=0.013\) & 0.846 & 0.846\\
Modified Hachinski score & ADAS-Cog 13
 & 212 & \(\rho=0.029\) & 0.674 & 0.846\\
Modified Hachinski score & CDR memory
 & 227 & \(\rho=0.043\) & 0.516 & 0.846\\
Modified Hachinski score & MoCA
 & 130 & \(\rho=-0.034\) & 0.704 & 0.846\\
Modified Hachinski score & Normalized WMH burden
 & 227 & \(\rho=0.004\) & 0.949 & 0.949\\
\addlinespace
Normalized WMH burden & Average controllability
 & 278 & \(\rho=-0.142\) & 0.018 & 0.044\\
Normalized WMH burden & Modal controllability
 & 278 & \(\rho=0.137\) & 0.022 & 0.044\\
Normalized WMH burden & Hippocampal strength
 & 278 & \(\rho=-0.110\) & 0.067 & 0.089\\
Normalized WMH burden & Hippocampal betweenness
 & 278 & \(\rho=-0.062\) & 0.304 & 0.304\\
Normalized WMH burden & MMSE
 & 276 & \(\rho=-0.122\) & 0.043 & 0.108\\
Normalized WMH burden & ADAS-Cog 11
 & 262 & \(\rho=0.082\) & 0.186 & 0.229\\
Normalized WMH burden & ADAS-Cog 13
 & 262 & \(\rho=0.075\) & 0.229 & 0.229\\
Normalized WMH burden & CDR memory
 & 277 & \(\rho=0.104\) & 0.084 & 0.140\\
Normalized WMH burden & MoCA
 & 167 & \(\rho=-0.247\) & 0.0013 & 0.006\\
\bottomrule
\end{tabularx}

\vspace{1mm}
\footnotesize{\textit{Note:} Modified Hachinski scores and UCD
white-matter hyperintensity (WMH) measurements were matched to the
selected diffusion MRI session within 365 days. WMH burden was calculated
as \(100\times\mathrm{TOTAL\_WMH}/\mathrm{CEREBRUM\_TCV}\) and transformed
using \(\log(1+x)\). Partial Spearman associations controlled for clinical
group, age, sex, and education. BH-FDR correction was applied separately
within each vascular predictor and outcome family. }
\end{table}

\begin{table}[!htbp]
\centering
\caption{Harmonized shared and incremental information analyses using exact-date FreeSurfer 7.2.0 hippocampal volume and intracranial volume.}
\label{tab:tableS6}
\scriptsize

\textbf{A. Variance in unranked controllability explained by volume and topology}

\vspace{1mm}
\begin{tabular}{@{}lccc@{}}
\toprule
\textbf{Outcome} & \textbf{$n$} & \textbf{Model $R^2$ (95\% CI)} & \textbf{Adjusted $R^2$}\\
\midrule
Raw average controllability & 270 & 0.546 (0.449--0.688) & 0.533 \\
Raw modal controllability   & 270 & 0.550 (0.481--0.655) & 0.536 \\
\bottomrule
\end{tabular}

\vspace{3mm}
\textbf{B. Increment after jointly adding average- and modal-controllability ranks}

\vspace{1mm}
\scriptsize
\setlength{\tabcolsep}{3pt}
\begin{tabular}{@{}lcccccccc@{}}
\toprule
\textbf{Outcome} &
\textbf{$n$} &
\textbf{Base $R^2$} &
\textbf{Full $R^2$} &
\textbf{$\Delta R^2$ (95\% CI)} &
\textbf{Partial $F$} &
\textbf{df} &
\textbf{$p$} &
\textbf{$p_{\mathrm{FDR}}$}\\
\midrule
MMSE        & 268 & 0.388 & 0.403 & 0.0149 (0.0000--0.0530) & 3.213 & 2,257 & 0.0418 & 0.0974 \\
ADAS-Cog 11 & 254 & 0.421 & 0.431 & 0.0097 (0.0000--0.0364) & 2.079 & 2,243 & 0.1273 & 0.1273 \\
ADAS-Cog 13 & 254 & 0.444 & 0.454 & 0.0101 (0.0000--0.0359) & 2.257 & 2,243 & 0.1068 & 0.1273 \\
CDR memory  & 269 & 0.332 & 0.357 & 0.0249 (0.0047--0.0638) & 4.991 & 2,258 & 0.0075 & 0.0374 \\
MoCA        & 166 & 0.322 & 0.347 & 0.0244 (0.0000--0.0758) & 2.893 & 2,155 & 0.0584 & 0.0974 \\
\bottomrule
\end{tabular}

\vspace{3mm}
\textbf{C. Increment for multinomial clinical-group membership}

\vspace{1mm}
\begin{tabular}{@{}lccccccc@{}}
\toprule
\textbf{Outcome} & \textbf{$n$} & \textbf{Base McFadden $R^2$} & \textbf{Full McFadden $R^2$} &
\textbf{$\Delta R^{2}_{\mathrm{McFadden}}$ (95\% CI)} & \textbf{$\chi^2$} & \textbf{df} & \textbf{$p$}\\
\midrule
Clinical-group membership & 270 & 0.1885 & 0.2105 & 0.0219 (0.0103--0.0606) & 15.180 & 6 & 0.0189 \\
\bottomrule
\end{tabular}

\vspace{3mm}
\textbf{D. Availability of exact-date FreeSurfer variables}\par\vspace{1mm}
\begin{tabular}{@{}lrrrrrrrr@{}}
\toprule
\textbf{Group} & \textbf{Cohort} & \textbf{FS complete} & \textbf{FS missing} & \textbf{MMSE} & \textbf{ADAS-11} & \textbf{ADAS-13} & \textbf{CDR memory} & \textbf{MoCA} \\
\midrule
CN & 110 & 105 & 5 & 104 & 104 & 104 & 105 & 77 \\
SMC & 27 & 27 & 0 & 27 & 27 & 27 & 27 & 3 \\
MCI & 88 & 85 & 3 & 84 & 73 & 73 & 84 & 67 \\
AD & 56 & 53 & 3 & 53 & 50 & 50 & 53 & 19 \\
Total & 281 & 270 & 11 & 268 & 254 & 254 & 269 & 166 \\
\bottomrule
\end{tabular}

\vspace{1mm}
\footnotesize{\textit{Note:} ICV, left hippocampal volume, and right hippocampal volume were \texttt{ST10CV}, \texttt{ST29SV}, and \texttt{ST88SV} from \texttt{UCSFFSX7\_11Jul2026.csv}. Eligible 3.0~T records exactly matched the diffusion-MRI date; lower numeric \texttt{IMAGEUID} resolved duplicates. All included records used FreeSurfer 7.2.0. \texttt{STATUS} was \texttt{partial}; \texttt{OVERALLQC} and \texttt{HIPPOQC} were not populated. Panel A base models included age, sex, education, ICV, left and right hippocampal volumes as separate predictors, strength, and betweenness. Panel B jointly added rank-based average and modal controllability and used 5,000 group-stratified percentile bootstrap resamples; the intervals are descriptive, whereas inference used the BH-FDR-corrected partial \(F\)-tests across five outcomes. Panel C used a multinomial likelihood-ratio test and a 5,000-resample group-stratified percentile-bootstrap interval.}

\end{table}

\begin{table}[!htbp]
\centering
\caption{Eddy-derived motion and motion-adjusted sensitivity analyses of hippocampal controllability.}
\label{tab:tableS7}
\small
\setlength{\tabcolsep}{4pt}
\renewcommand{\arraystretch}{1.12}
\begin{tabularx}{\textwidth}{@{}>{\raggedright\arraybackslash}X
                                  >{\raggedright\arraybackslash}p{3.0cm}
                                  cccc@{}}
\toprule
\textbf{Adjustment model or test} &
\textbf{Measure} &
\textbf{\(n\)} &
\textbf{Statistic} &
\textbf{\(p\) value} &
\textbf{\(p_{\mathrm{FDR}}\)}\\
\midrule
\multicolumn{6}{l}{\textbf{A. Clinical-group differences in eddy-derived motion}}\\
Kruskal--Wallis group test & Mean absolute RMS    & 281 & \(H=4.347\) & 0.226 & 0.285\\
Kruskal--Wallis group test & Mean relative RMS    & 281 & \(H=6.448\) & 0.092 & 0.233\\
Kruskal--Wallis group test & Maximum absolute RMS & 281 & \(H=3.787\) & 0.285 & 0.285\\
Kruskal--Wallis group test & Maximum relative RMS & 281 & \(H=5.903\) & 0.116 & 0.233\\
\addlinespace
\multicolumn{6}{l}{\textbf{B. Sensitivity analyses after motion adjustment}}\\
Demographic covariates
  & Average controllability & 281 & \(F=6.495\) & 0.0004 & 0.0007\\
  & Modal controllability   & 281 & \(F=6.041\) & 0.0007 & 0.0007\\
\addlinespace
Demographic covariates + motion
  & Average controllability & 281 & \(F=6.237\) & 0.0006 & 0.0012\\
  & Modal controllability   & 281 & \(F=5.786\) & 0.0014 & 0.0014\\
\bottomrule
\end{tabularx}

\vspace{1mm}
\footnotesize{\textit{Note:} RMS displacement measures were obtained from FSL \texttt{eddy} outputs. Panel A reports Kruskal--Wallis tests, with BH-FDR correction across the four motion summaries. Panel B reports Freedman--Lane permutation \(F\)-tests with 10,000 permutations. All Panel B models included age, sex, and education; the motion-adjusted model additionally included mean absolute and mean relative RMS displacement. BH-FDR correction was applied across average and modal controllability within each Panel B model. RMS, root mean square.}
\end{table}

\begin{table}[!htbp]
\centering
\caption{Multicenter and acquisition-protocol sensitivity analyses of hippocampal controllability.}
\label{tab:tableS8}
\small
\setlength{\tabcolsep}{4pt}
\renewcommand{\arraystretch}{1.12}
\begin{tabularx}{\textwidth}{@{}>{\raggedright\arraybackslash}X
                                  >{\raggedright\arraybackslash}p{3.0cm}
                                  cccc@{}}
\toprule
\textbf{Adjustment model} &
\textbf{Measure} &
\textbf{\(n\)} &
\textbf{\(F\)} &
\textbf{\(p_{\mathrm{perm}}\)} &
\textbf{\(p_{\mathrm{FDR}}\)}\\
\midrule
Scanner manufacturer + ADNI phase
  & Average controllability & 281 & 5.362 & 0.0022 & 0.0026\\
  & Modal controllability   & 281 & 4.852 & 0.0026 & 0.0026\\
\addlinespace
Scanner manufacturer + phase + site
  & Average controllability & 281 & 3.792 & 0.0119 & 0.0166\\
  & Modal controllability   & 281 & 3.522 & 0.0166 & 0.0166\\
\addlinespace
Scanner/site/protocol/motion model
  & Average controllability & 281 & 4.368 & 0.0057 & 0.0093\\
  & Modal controllability   & 281 & 3.977 & 0.0093 & 0.0093\\
\bottomrule
\end{tabularx}

\vspace{1mm}
\footnotesize{\textit{Note:} All models included age, sex, and education and were evaluated using Freedman--Lane permutation \(F\)-tests with 10,000 permutations. After reconciliation to the final NIfTI inputs, the strict scanner/site/protocol/motion model additionally included 12 scanner-model levels, three ADNI phases, 43 sites, two acceleration levels, nine NIfTI-derived diffusion-volume-count levels, four slice-thickness levels, and mean absolute and relative RMS displacement. Its nominal 76-column full design had rank 71 and 210 residual degrees of freedom. BH-FDR correction was applied across average and modal controllability within each adjustment model. ComBat was not applied because diagnosis, ADNI phase, scanner, and site were not fully crossed.}
\end{table}

\begin{table}[!htbp]
\centering
\caption{SMC sample-size and unequal-group sensitivity analyses of hippocampal controllability.}
\label{tab:tableS9}
\scriptsize
\setlength{\tabcolsep}{3pt}
\renewcommand{\arraystretch}{1.12}
\begin{tabularx}{\textwidth}{@{}>{\raggedright\arraybackslash}p{2.4cm}
                                  >{\raggedright\arraybackslash}p{1.7cm}
                                  >{\raggedright\arraybackslash}X
                                  >{\raggedright\arraybackslash}X
                                  >{\centering\arraybackslash}p{1.7cm}
                                  >{\centering\arraybackslash}p{1.7cm}@{}}
\toprule
\textbf{Analysis} &
\textbf{Measure} &
\textbf{Partial \(F\)} &
\textbf{\(\Delta R^2\)} &
\textbf{Significant runs} &
\textbf{SMC pattern retained}\\
\midrule
Balanced subsampling
  & Average & 5.663 (3.439--9.859) & 0.131 (0.084--0.212) & 99.9\% & 99.9\%\\
  & Modal   & 5.349 (3.295--9.588) & 0.127 (0.081--0.205) & 99.7\% & 99.9\%\\
\addlinespace
Leave-one-SMC-out
  & Average & 5.960--6.863 & 0.058--0.066 & 100\% & 100\%\\
  & Modal   & 5.518--6.398 & 0.054--0.062 & 100\% & 100\%\\
\bottomrule
\end{tabularx}

\vspace{1mm}
\footnotesize{\textit{Note:} Balanced subsampling comprised 1,000 repetitions with 27 participants per group; values are medians (2.5th--97.5th percentiles), and ``significant runs'' denotes the percentage with parametric partial \(F\)-test \(p<0.05\). Leave-one-SMC-out comprised 27 analyses with 26 SMC participants and all CN, MCI, and AD participants retained; values are ranges across iterations, all of which remained significant by Freedman--Lane tests with 10,000 permutations (average controllability, \(p=0.0002\text{--}0.0010\); modal controllability, \(p=0.0003\text{--}0.0017\)). The SMC pattern was defined as the highest adjusted mean average controllability and lowest adjusted mean modal controllability among the four groups. Overlapping resamples are stability checks rather than independent replications.}
\end{table}

\begin{table}[!htbp]
\centering
\caption{Same-day FreeSurfer global gray-matter sensitivity analysis.}
\label{tab:tableS10}
\footnotesize

\textbf{A. Same-day coverage}

\vspace{1mm}
\begin{tabular}{@{}lccc@{}}
\toprule
\textbf{Group} & \textbf{Cohort \(n\)} & \textbf{Matched \(n\)} & \textbf{Coverage}\\
\midrule
CN & 110 & 105 & 95.5\% \\
SMC & 27 & 27 & 100.0\% \\
MCI & 88 & 85 & 96.6\% \\
AD & 56 & 53 & 94.6\% \\
Overall & 281 & 270 & 96.1\% \\
\bottomrule
\end{tabular}

\vspace{3mm}
\textbf{B. Partial Spearman associations}

\vspace{1mm}
\setlength{\tabcolsep}{3pt}
\begin{tabular}{@{}>{\raggedright\arraybackslash}p{4.0cm}>{\raggedright\arraybackslash}p{4.3cm}ccc@{}}
\toprule
\textbf{Hippocampal measure} & \textbf{Comparison measure} & \textbf{\(n\)} & \textbf{\(\rho\)} & \textbf{\(p_{\mathrm{FDR}}\)}\\
\midrule
Bilateral hippocampus & Total gray matter & 270 & 0.498 & \(7.54\times10^{-18}\) \\
Bilateral hippocampus & Non-hippocampal gray matter & 270 & 0.480 & \(1.06\times10^{-16}\) \\
Bilateral hippocampus & Cortical gray matter & 270 & 0.468 & \(5.01\times10^{-16}\) \\
Bilateral hippocampus & Subcortical gray matter & 270 & 0.595 & \(1.91\times10^{-26}\) \\
Left hippocampus & Non-hippocampal gray matter & 270 & 0.470 & \(4.48\times10^{-16}\) \\
Right hippocampus & Non-hippocampal gray matter & 270 & 0.440 & \(3.05\times10^{-14}\) \\
\bottomrule
\end{tabular}

\vspace{3mm}
\textbf{C. Associations additionally adjusted for diagnostic group}

\vspace{1mm}
\begin{tabular}{@{}>{\raggedright\arraybackslash}p{4.0cm}>{\raggedright\arraybackslash}p{4.3cm}ccc@{}}
\toprule
\textbf{Hippocampal measure} & \textbf{Comparison measure} & \textbf{\(n\)} & \textbf{\(\rho\)} & \textbf{\(p_{\mathrm{FDR}}\)}\\
\midrule
Bilateral hippocampus & Total gray matter & 270 & 0.428 & \(5.67\times10^{-13}\) \\
Bilateral hippocampus & Non-hippocampal gray matter & 270 & 0.412 & \(3.38\times10^{-12}\) \\
Bilateral hippocampus & Cortical gray matter & 270 & 0.395 & \(2.06\times10^{-11}\) \\
Bilateral hippocampus & Subcortical gray matter & 270 & 0.537 & \(7.86\times10^{-21}\) \\
Left hippocampus & Non-hippocampal gray matter & 270 & 0.404 & \(7.63\times10^{-12}\) \\
Right hippocampus & Non-hippocampal gray matter & 270 & 0.366 & \(5.52\times10^{-10}\) \\
\bottomrule
\end{tabular}

\vspace{1mm}
\footnotesize{\textit{Note:} Panel B associations controlled for ICV, age, sex, and education. Panel C additionally included three diagnostic-group indicator variables. Non-hippocampal gray matter was calculated by subtracting bilateral hippocampal volume from total gray-matter volume, avoiding part--whole correlation. BH-FDR correction was applied across the six associations within each panel. These are cross-sectional volume associations, not longitudinal atrophy rates. The matched FreeSurfer records lacked final visual QC and these results are exploratory.}

\end{table}

\begin{table}[!htbp]
\centering
\caption{Amyloid PET sensitivity analyses for hippocampal controllability.}
\label{tab:tableS11}
\scriptsize
\setlength{\tabcolsep}{3pt}
\begin{tabular}{@{}cllcccc@{}}
\toprule
\textbf{Window} & \textbf{Amyloid variable} & \textbf{Tested effect} & \textbf{Outcome} & \textbf{$n$} & \textbf{$F$} & \textbf{$p$ ($p_{\mathrm{FDR}}$)}\\
\midrule
365 d & Status & Diagnosis & Average & 230 & $F_{3,222}=5.521$ & 0.0018 (0.0018) \\
365 d & Status & Diagnosis & Modal & 230 & $F_{3,222}=5.321$ & 0.0016 (0.0018) \\
365 d & Status & Amyloid & Average & 230 & $F_{1,222}=0.327$ & 0.5689 (0.5689) \\
365 d & Status & Amyloid & Modal & 230 & $F_{1,222}=0.392$ & 0.5381 (0.5689) \\
365 d & Centiloids & Diagnosis & Average & 226 & $F_{3,218}=5.483$ & 0.0009 (0.0018) \\
365 d & Centiloids & Diagnosis & Modal & 226 & $F_{3,218}=5.160$ & 0.0021 (0.0021) \\
365 d & Centiloids & Amyloid & Average & 226 & $F_{1,218}=0.877$ & 0.3463 (0.3675) \\
365 d & Centiloids & Amyloid & Modal & 226 & $F_{1,218}=0.781$ & 0.3675 (0.3675) \\
\midrule
730 d & Status & Diagnosis & Average & 241 & $F_{3,233}=6.228$ & 0.0004 (0.0008) \\
730 d & Status & Diagnosis & Modal & 241 & $F_{3,233}=6.100$ & 0.0008 (0.0008) \\
730 d & Status & Amyloid & Average & 241 & $F_{1,233}=0.023$ & 0.8797 (0.8797) \\
730 d & Status & Amyloid & Modal & 241 & $F_{1,233}=0.051$ & 0.8246 (0.8797) \\
730 d & Centiloids & Diagnosis & Average & 237 & $F_{3,229}=6.106$ & 0.0006 (0.0008) \\
730 d & Centiloids & Diagnosis & Modal & 237 & $F_{3,229}=5.850$ & 0.0008 (0.0008) \\
730 d & Centiloids & Amyloid & Average & 237 & $F_{1,229}=0.394$ & 0.5188 (0.5628) \\
730 d & Centiloids & Amyloid & Modal & 237 & $F_{1,229}=0.336$ & 0.5628 (0.5628) \\
\bottomrule
\end{tabular}

\vspace{1mm}
\footnotesize{\textit{Note:} The diagnosis effect was tested after adjustment for the listed amyloid variable, age, sex, and education. The amyloid effect was tested after adjustment for diagnosis, age, sex, and education. Status and Centiloids were analyzed in separate complete-case models, producing different sample sizes; valid negative Centiloid values were retained. All \(p\) values were obtained using 10,000 Freedman--Lane permutations. BH-FDR correction was applied across average and modal controllability within each tested effect.}

\end{table}


\begin{table}[!htbp]
\centering
\caption{Classification performance based on hippocampal structural connectivity features in four binary tasks.}
\label{tab:tableS12}
\footnotesize
\setlength{\tabcolsep}{4pt}
\begin{tabular}{@{}lcc@{}}
\toprule
\textbf{Task} & \textbf{AUC (95\% CI)} & \textbf{Interpretation}\\
\midrule
CN vs AD & 0.601 (0.508--0.691) & Descriptive exploratory estimate \\
SMC vs MCI & 0.693 (0.584--0.797) & Descriptive exploratory estimate \\
MCI vs AD & 0.634 (0.540--0.731) & Descriptive exploratory estimate \\
CN vs SMC & 0.680 (0.573--0.780) & Descriptive exploratory estimate \\

\bottomrule
\end{tabular}

\vspace{1mm}
\footnotesize{\textit{Note:} AUC confidence intervals were estimated using 2,000 bootstrap resamples of subject-level out-of-sample predictions from nested cross-validation. The previous fixed-score permutation values were removed because the full high-dimensional pipeline was not refitted under label permutation. CN, cognitively normal; SMC, subjective memory complaints; MCI, mild cognitive impairment; AD, Alzheimer's disease.}

\end{table}

% Begin inlined supplementary_new_tables.tex
\begin{table}[!htbp]
\centering
\caption{Primary rank endpoints and matrix-scale-adjusted unranked hippocampal controllability sensitivity analyses.}
\label{tab:rank_raw_sensitivity}
\scriptsize
\setlength{\tabcolsep}{3.2pt}
\textbf{A. Adjusted group estimates (95\% CI)}\par\vspace{1mm}
\begin{tabular}{llcccc}
\toprule
Scale & Metric & CN & SMC & MCI & AD\\
\midrule
Within-subject rank & Average & 345.3 (329.5--360.3) & 383.9 (366.1--399.6) & 324.7 (306.7--341.5) & 307.4 (286.7--327.7)\\
Within-subject rank & Modal & 107.9 (92.5--124.0) & 66.5 (49.4--85.0) & 128.8 (111.5--148.0) & 142.4 (122.5--163.9)\\
Unranked raw value & Average & 1.173 (1.152--1.195) & 1.218 (1.174--1.265) & 1.169 (1.140--1.204) & 1.116 (1.096--1.139)\\
Unranked raw value & Modal & 0.905 (0.894--0.914) & 0.877 (0.858--0.895) & 0.913 (0.901--0.924) & 0.929 (0.918--0.940)\\
\bottomrule
\end{tabular}

\vspace{3mm}
\textbf{B. Omnibus clinical-group effects}\par\vspace{1mm}
\begin{tabular}{llccccc}
\toprule
Analysis & Metric & $F$ & df & $p_{\mathrm{perm}}$ & $p_{\mathrm{FDR}}$ & Partial $\eta^2$\\
\midrule
Primary rank & Average & 6.495 & 3,274 & 0.0004 & 0.0028 & 0.066 (0.034--0.127)\\
Primary rank & Modal & 6.041 & 3,274 & 0.0007 & 0.0049 & 0.062 (0.032--0.121)\\
Raw, matrix-scale adjusted & Average & 4.426 & 3,273 & 0.0066 & 0.0066 & 0.046\\
Raw, matrix-scale adjusted & Modal & 6.313 & 3,273 & 0.0004 & 0.0008 & 0.065\\
Raw, matrix-scale+density adjusted & Average & 4.232 & 3,272 & 0.0062 & 0.0062 & 0.045\\
Raw, matrix-scale+density adjusted & Modal & 5.306 & 3,272 & 0.0016 & 0.0032 & 0.055\\
\bottomrule
\end{tabular}

\vspace{3mm}
\textbf{C. Direct average--modal association}\par\vspace{1mm}
\begin{tabular}{lcc}
\toprule
Scale and adjustment & $n$ & Partial Spearman $\rho$ (95\% CI)\\
\midrule
Within-subject ranks; age, sex, education & 281 & $-0.992$ ($-0.994$ to $-0.987$)\\
Unranked; age, sex, education, matrix scale & 281 & $-0.930$ ($-0.955$ to $-0.896$)\\
\bottomrule
\end{tabular}
\par\vspace{1mm}
\footnotesize{\textit{Note:} Adjusted estimates and confidence intervals used 5,000 group-stratified bootstrap resamples. Omnibus tests used 10,000 Freedman--Lane permutations. Primary rank values retain correction across seven Tian structures within endpoint; each unranked omnibus family contains average and modal outcomes. Matrix scale is $\log(1+W_{\mathrm{total}})$, where $W_{\mathrm{total}}=\sum_{i<j}A_{ij}$ for the final inverse-node-volume-scaled matrix before spectral stabilization.}
\end{table}

\begin{table}[!htbp]
\centering
\caption{Multiplicity-controlled pairwise contrasts for the primary hippocampal within-subject rank endpoints.}
\label{tab:primary_pairwise_contrasts}
\scriptsize
\begin{tabular}{lllrrr}
\toprule
Endpoint & Contrast & Adjusted difference (95\% CI) & Standardized difference & $p$ & $p_{\mathrm{FDR}}$\\
\midrule
Average rank & SMC--CN & 38.6 (15.5, 60.9) & 0.49 & 0.0225 & 0.0337\\
& MCI--CN & $-20.6$ ($-43.7$, 2.0) & $-0.26$ & 0.0733 & 0.0977\\
& AD--CN & $-37.9$ ($-63.9$, $-12.4$) & $-0.48$ & 0.0049 & 0.0118\\
& MCI--SMC & $-59.2$ ($-83.3$, $-34.3$) & $-0.75$ & 0.0011 & 0.0033\\
& AD--SMC & $-76.5$ ($-102.6$, $-50.0$) & $-0.97$ & 0.0002 & 0.0018\\
& AD--MCI & $-17.3$ ($-43.6$, 8.8) & $-0.22$ & 0.2027 & 0.2211\\
\addlinespace
Modal rank & SMC--CN & $-41.5$ ($-65.1$, $-17.2$) & $-0.51$ & 0.0196 & 0.0336\\
& MCI--CN & 20.9 ($-2.3$, 45.6) & 0.26 & 0.0823 & 0.0988\\
& AD--CN & 34.5 (8.8, 61.3) & 0.43 & 0.0126 & 0.0252\\
& MCI--SMC & 62.3 (36.4, 88.2) & 0.77 & 0.0007 & 0.0028\\
& AD--SMC & 76.0 (48.5, 102.3) & 0.94 & 0.0003 & 0.0018\\
& AD--MCI & 13.6 ($-13.5$, 40.5) & 0.17 & 0.3317 & 0.3317\\
\bottomrule
\end{tabular}
\par\vspace{1mm}
\footnotesize{\textit{Note:} Differences are first-named minus second-named group. Confidence intervals used 5,000 group-stratified bootstrap resamples and are individually unadjusted. The 12 two-sided permutation tests formed one BH-FDR family.}
\end{table}

\begin{table}[!htbp]
\centering
\caption{Multiplicity-controlled pairwise contrasts for matrix-scale-adjusted unranked hippocampal controllability.}
\label{tab:raw_pairwise}
\scriptsize
\begin{tabular}{lllrrr}
\toprule
Metric & Contrast & Adjusted difference (95\% CI) & Standardized difference & $p_{\mathrm{FDR}}$\\
\midrule
Average & SMC--CN & 0.0446 ($-0.0059$, 0.0966) & 0.36 & 0.1193\\
& MCI--CN & $-0.0040$ ($-0.0408$, 0.0350) & $-0.03$ & 0.8259\\
& AD--CN & $-0.0569$ ($-0.0879$, $-0.0244$) & $-0.46$ & 0.0187\\
& MCI--SMC & $-0.0486$ ($-0.1032$, 0.0078) & $-0.39$ & 0.1108\\
& AD--SMC & $-0.1016$ ($-0.1532$, $-0.0529$) & $-0.82$ & 0.0048\\
& AD--MCI & $-0.0529$ ($-0.0947$, $-0.0148$) & $-0.43$ & 0.0255\\
\addlinespace
Modal & SMC--CN & $-0.0272$ ($-0.0485$, $-0.0058$) & $-0.53$ & 0.0255\\
& MCI--CN & 0.0083 ($-0.0073$, 0.0236) & 0.16 & 0.2775\\
& AD--CN & 0.0245 (0.0092, 0.0404) & 0.48 & 0.0135\\
& MCI--SMC & 0.0356 (0.0133, 0.0579) & 0.70 & 0.0092\\
& AD--SMC & 0.0517 (0.0294, 0.0742) & 1.01 & 0.0012\\
& AD--MCI & 0.0161 (0.0003, 0.0327) & 0.32 & 0.1012\\
\bottomrule
\end{tabular}
\par\vspace{1mm}
\footnotesize{\textit{Note:} Models additionally adjusted for $\log(1+W_{\mathrm{total}})$, where $W_{\mathrm{total}}=\sum_{i<j}A_{ij}$ for the final inverse-node-volume-scaled matrix before spectral stabilization. The 12 two-sided permutation tests formed one BH-FDR family; individual bootstrap confidence intervals are unadjusted.}
\end{table}

\begin{landscape}
\begin{longtable}{llrrrrr}
\caption{Direct paired AUC comparisons of candidate feature sets with bilateral hippocampal volume in common complete-case samples.}\label{tab:paired_volume_classification}\\
\toprule
Task & Candidate & $n$ & Candidate AUC & Volume AUC & $\Delta$AUC (95\% CI) & $p_{\mathrm{Holm}}$\\
\midrule
\endfirsthead
\toprule Task & Candidate & $n$ & Candidate AUC & Volume AUC & $\Delta$AUC (95\% CI) & $p_{\mathrm{Holm}}$\\ \midrule
\endhead
CN vs AD & Global controllability & 158 & 0.653 & 0.820 & $-0.167$ ($-0.281$, $-0.057$) & 0.042\\
& Left controllability & 158 & 0.673 & 0.820 & $-0.147$ ($-0.256$, $-0.039$) & 0.120\\
& Right controllability & 158 & 0.587 & 0.820 & $-0.233$ ($-0.344$, $-0.120$) & 0.007\\
& Topology & 158 & 0.549 & 0.820 & $-0.271$ ($-0.384$, $-0.148$) & 0.004\\
& Combined & 158 & 0.813 & 0.820 & $-0.007$ ($-0.037$, 0.025) & 1.000\\
\addlinespace
SMC vs MCI & Global controllability & 112 & 0.708 & 0.654 & 0.054 ($-0.091$, 0.199) & 1.000\\
& Left controllability & 112 & 0.675 & 0.654 & 0.022 ($-0.122$, 0.167) & 1.000\\
& Right controllability & 112 & 0.708 & 0.654 & 0.055 ($-0.098$, 0.200) & 1.000\\
& Topology & 112 & 0.617 & 0.654 & $-0.037$ ($-0.196$, 0.119) & 1.000\\
& Combined & 112 & 0.729 & 0.654 & 0.076 ($-0.014$, 0.169) & 1.000\\
\addlinespace
MCI vs AD & Global controllability & 138 & 0.630 & 0.825 & $-0.195$ ($-0.319$, $-0.067$) & 0.019\\
& Left controllability & 138 & 0.649 & 0.825 & $-0.176$ ($-0.292$, $-0.058$) & 0.033\\
& Right controllability & 138 & 0.454 & 0.825 & $-0.370$ ($-0.499$, $-0.238$) & 0.004\\
& Topology & 138 & 0.490 & 0.825 & $-0.335$ ($-0.456$, $-0.212$) & 0.004\\
& Combined & 138 & 0.808 & 0.825 & $-0.016$ ($-0.042$, 0.006) & 1.000\\
\addlinespace
CN vs SMC & Global controllability & 132 & 0.606 & 0.635 & $-0.029$ ($-0.194$, 0.138) & 1.000\\
& Left controllability & 132 & 0.577 & 0.635 & $-0.057$ ($-0.214$, 0.105) & 1.000\\
& Right controllability & 132 & 0.632 & 0.635 & $-0.003$ ($-0.165$, 0.163) & 1.000\\
& Topology & 132 & 0.672 & 0.635 & 0.037 ($-0.119$, 0.199) & 1.000\\
& Combined & 132 & 0.697 & 0.635 & 0.062 ($-0.028$, 0.153) & 1.000\\
\bottomrule
\end{longtable}
\footnotesize\textit{Note:} Candidate and volume models used identical subjects, cross-validation splits, and the same random seed as Table~II. Confidence intervals used 5,000 group-stratified resamples of fixed averaged OOF scores and therefore do not include model-refitting variability. Paired score-swap tests were Holm corrected across 20 comparisons.
\end{landscape}

\begin{table}[!htbp]
\centering
\caption{Reconstructed query inventory and partial cohort-selection audit from the ADNI exports and analysis archive.}
\label{tab:cohort_flow_audit}
\footnotesize
\begin{tabular}{p{0.29\textwidth}rrp{0.36\textwidth}}
\toprule
Stage & Individuals & Sessions/series & Rule or audit status\\
\midrule
MRIQC export: all imaging-series rows & 3,194 & 92,098 & Imaging-series export used for analysis\\
3.0~T diffusion-MRI query & 1,860 & 5,793 & dMRI and 3.0~T\\
Unique participant-date sessions & 1,860 & 4,566 & Same-date series collapsed; lower image ID broke ties\\
Eligibility and deterministic session selection & -- & -- & Complete gradients/phase encoding, matched T1, usable QC, successful registration/tractography/connectome; earliest eligible date\\
Previous revised cohort & 240 & 240 & CN 110, SMC 18, MCI 88, AD 24; all retained\\
Additional eligible participants & 41 & 41 & SMC 9 and AD 32; same eligibility criteria\\
Final primary cohort & 281 & 281 & One selected session per participant; CN 110, SMC 27, MCI 88, AD 56\\
\bottomrule
\end{tabular}
\par\vspace{1mm}
\footnotesize{\textit{Note:} A session was a unique participant-date combination. Same-date diffusion series were resolved using the lower numeric image identifier; when several eligible dates remained, the earliest eligible date was selected. The final cohort retained the previous 240 participants and added 41 eligible participants identified during the re-audit. The query and final session inventory were recoverable, whereas intermediate historical exclusion counts and reasons were not fully preserved; the table is therefore labeled as a partial audit rather than a complete prospective flow.}
\end{table}

% Begin inlined acquisition_cells_R2.tex
% Generated from the corrected final NIfTI input manifest.
\begin{landscape}
\scriptsize
\setlength{\tabcolsep}{2.5pt}
\begin{longtable}{lllllrrrrrrr}
\caption{Cross-classification of clinical group by ADNI phase, site, scanner, and principal diffusion-protocol characteristics.}\label{tab:acquisition_cells} \\
\toprule
Phase & Site & Manufacturer & Model & Acceleration & NIfTI volumes & Slice (mm) & CN & SMC & MCI & AD & Total \\
\midrule
\endfirsthead
\multicolumn{12}{c}{\tablename\ \thetable{} -- continued} \\
\toprule
Phase & Site & Manufacturer & Model & Acceleration & NIfTI volumes & Slice (mm) & CN & SMC & MCI & AD & Total \\
\midrule
\endhead
\midrule
\multicolumn{12}{r}{Continued on next page} \\
\endfoot
\bottomrule
\endlastfoot
ADNI3 & 168 & Siemens & Prisma\_fit & Accelerated & 55 & 2.0 & 5 & 0 & 7 & 7 & 19 \\
ADNI3 & 135 & GE & DISCOVERY MR750w & Accelerated & 36 & 2.0 & 4 & 0 & 8 & 4 & 16 \\
ADNIGO/ADNI2 & 127 & GE & DISCOVERY MR750 & Accelerated & 46 & 2.7 & 2 & 3 & 0 & 6 & 11 \\
ADNIGO/ADNI2 & 27 & GE & DISCOVERY MR750 & Accelerated & 46 & 2.7 & 2 & 6 & 0 & 1 & 9 \\
ADNI3 & 3 & Siemens & Prisma & Accelerated & 55 & 2.0 & 4 & 0 & 5 & 0 & 9 \\
ADNI3 & 35 & Siemens & Prisma\_fit & Accelerated & 55 & 2.0 & 4 & 0 & 2 & 1 & 7 \\
ADNI3 & 127 & GE & DISCOVERY MR750 & Accelerated & 54 & 2.0 & 4 & 0 & 1 & 2 & 7 \\
ADNI3 & 141 & Siemens & Verio & Unaccelerated & 31 & 2.0 & 5 & 0 & 2 & 0 & 7 \\
ADNI3 & 109 & Siemens & TrioTim & Accelerated & 31 & 2.0 & 0 & 0 & 7 & 0 & 7 \\
ADNIGO/ADNI2 & 16 & GE & Signa HDxt & Accelerated & 46 & 2.7 & 0 & 0 & 0 & 7 & 7 \\
ADNI3 & 116 & Siemens & TrioTim & Accelerated & 31 & 2.0 & 4 & 0 & 1 & 1 & 6 \\
ADNIGO/ADNI2 & 3 & GE & Signa HDxt & Accelerated & 55 & 2.7 & 3 & 1 & 2 & 0 & 6 \\
ADNI3 & 14 & Siemens & Verio & Accelerated & 31 & 2.0 & 4 & 0 & 2 & 0 & 6 \\
ADNI3 & 37 & Siemens & Prisma & Accelerated & 55 & 2.0 & 3 & 0 & 3 & 0 & 6 \\
ADNI3 & 41 & Siemens & Prisma\_fit & Accelerated & 55 & 2.0 & 3 & 0 & 3 & 0 & 6 \\
ADNI3 & 129 & GE & DISCOVERY MR750 & Accelerated & 54 & 2.0 & 2 & 0 & 2 & 1 & 5 \\
ADNIGO/ADNI2 & 3 & GE & Signa HDxt & Accelerated & 46 & 2.7 & 2 & 1 & 0 & 2 & 5 \\
ADNI3 & 82 & Siemens & Verio & Accelerated & 31 & 2.0 & 4 & 0 & 1 & 0 & 5 \\
ADNI3 & 13 & Siemens & Biograph\_mMR & Accelerated & 31 & 2.0 & 1 & 0 & 2 & 1 & 4 \\
ADNIGO/ADNI2 & 29 & GE & Signa HDxt & Accelerated & 46 & 2.7 & 1 & 2 & 0 & 1 & 4 \\
ADNI3 & 11 & Siemens & Prisma\_fit & Accelerated & 55 & 2.0 & 3 & 0 & 1 & 0 & 4 \\
ADNI3 & 12 & Siemens & Prisma & Accelerated & 55 & 2.0 & 0 & 2 & 2 & 0 & 4 \\
ADNI3 & 27 & GE & DISCOVERY MR750 & Accelerated & 54 & 2.0 & 1 & 0 & 3 & 0 & 4 \\
ADNI3 & 68 & Siemens & Prisma\_fit & Accelerated & 55 & 2.0 & 3 & 0 & 1 & 0 & 4 \\
ADNI3 & 70 & Siemens & TrioTim & Accelerated & 31 & 2.0 & 2 & 0 & 2 & 0 & 4 \\
ADNIGO/ADNI2 & 7 & GE & Signa HDxt & Accelerated & 127 & 2.7 & 3 & 0 & 1 & 0 & 4 \\
ADNIGO/ADNI2 & 98 & GE & DISCOVERY MR750 & Accelerated & 46 & 2.7 & 2 & 0 & 0 & 2 & 4 \\
ADNI3 & 137 & Siemens & Skyra & Accelerated & 55 & 2.0 & 0 & 0 & 4 & 0 & 4 \\
ADNI3 & 9 & GE & DISCOVERY MR750w & Accelerated & 36 & 2.0 & 2 & 0 & 1 & 0 & 3 \\
ADNI3 & 16 & Siemens & Skyra & Accelerated & 55 & 2.0 & 2 & 0 & 1 & 0 & 3 \\
ADNI3 & 20 & Siemens & Prisma\_fit & Accelerated & 55 & 2.0 & 2 & 1 & 0 & 0 & 3 \\
ADNI3 & 22 & Siemens & TrioTim & Accelerated & 31 & 2.0 & 0 & 0 & 2 & 1 & 3 \\
ADNI3 & 33 & Siemens & Skyra & Accelerated & 55 & 2.0 & 2 & 0 & 1 & 0 & 3 \\
ADNI3 & 98 & GE & SIGNA Premier & Accelerated & 54 & 2.0 & 1 & 0 & 2 & 0 & 3 \\
ADNIGO/ADNI2 & 21 & GE & DISCOVERY MR750 & Accelerated & 68 & 2.7 & 1 & 2 & 0 & 0 & 3 \\
ADNIGO/ADNI2 & 27 & GE & DISCOVERY MR750 & Accelerated & 54 & 2.7 & 2 & 1 & 0 & 0 & 3 \\
ADNIGO/ADNI2 & 126 & GE & DISCOVERY MR750 & Accelerated & 46 & 2.7 & 1 & 0 & 0 & 2 & 3 \\
ADNIGO/ADNI2 & 127 & GE & DISCOVERY MR750 & Accelerated & 68 & 2.7 & 1 & 2 & 0 & 0 & 3 \\
ADNI3 & 32 & Siemens & Prisma\_fit & Accelerated & 55 & 2.0 & 3 & 0 & 0 & 0 & 3 \\
ADNIGO/ADNI2 & 5 & GE & DISCOVERY MR750 & Accelerated & 46 & 2.7 & 0 & 0 & 0 & 3 & 3 \\
ADNIGO/ADNI2 & 7 & GE & Signa HDxt & Accelerated & 46 & 2.7 & 0 & 0 & 0 & 3 & 3 \\
ADNI3 & 2 & Siemens & Prisma\_fit & Accelerated & 55 & 2.0 & 1 & 0 & 1 & 0 & 2 \\
ADNI3 & 24 & Siemens & Prisma\_fit & Accelerated & 127 & 2.0 & 1 & 0 & 1 & 0 & 2 \\
ADNI3 & 37 & Siemens & Prisma & Accelerated & 127 & 2.0 & 1 & 0 & 1 & 0 & 2 \\
ADNI3 & 114 & Siemens & Verio & Unaccelerated & 31 & 2.0 & 1 & 0 & 1 & 0 & 2 \\
ADNI3 & 126 & GE & DISCOVERY MR750 & Accelerated & 54 & 2.0 & 1 & 0 & 0 & 1 & 2 \\
ADNIGO/ADNI2 & 21 & GE & DISCOVERY MR750 & Accelerated & 54 & 2.7 & 1 & 1 & 0 & 0 & 2 \\
ADNIGO/ADNI2 & 27 & GE & Signa HDxt & Accelerated & 46 & 2.7 & 0 & 0 & 1 & 1 & 2 \\
ADNIGO/ADNI2 & 94 & GE & Signa HDxt & Accelerated & 46 & 2.7 & 1 & 0 & 0 & 1 & 2 \\
ADNIGO/ADNI2 & 99 & GE & DISCOVERY MR750 & Accelerated & 54 & 2.0 & 1 & 0 & 1 & 0 & 2 \\
ADNIGO/ADNI2 & 126 & GE & DISCOVERY MR750 & Accelerated & 68 & 2.7 & 0 & 1 & 1 & 0 & 2 \\
ADNIGO/ADNI2 & 127 & GE & DISCOVERY MR750 & Accelerated & 54 & 2.7 & 0 & 1 & 1 & 0 & 2 \\
ADNI3 & 21 & GE & DISCOVERY MR750 & Accelerated & 54 & 2.0 & 2 & 0 & 0 & 0 & 2 \\
ADNIGO/ADNI2 & 27 & GE & DISCOVERY MR750w & Accelerated & 46 & 2.7 & 0 & 0 & 0 & 2 & 2 \\
ADNIGO/ADNI2 & 57 & GE & DISCOVERY MR750 & Accelerated & 46 & 2.7 & 0 & 2 & 0 & 0 & 2 \\
ADNIGO/ADNI2 & 129 & GE & SIGNA HDx & Accelerated & 54 & 2.7 & 2 & 0 & 0 & 0 & 2 \\
ADNI3 & 2 & Siemens & Prisma\_fit & Accelerated & 127 & 2.0 & 0 & 0 & 1 & 0 & 1 \\
ADNI3 & 5 & GE & DISCOVERY MR750 & Accelerated & 54 & 2.0 & 0 & 0 & 1 & 0 & 1 \\
ADNI3 & 21 & GE & SIGNA UHP & Accelerated & 36 & 2.0 & 1 & 0 & 0 & 0 & 1 \\
ADNI3 & 29 & GE & DISCOVERY MR750w & Accelerated & 36 & 2.0 & 0 & 0 & 1 & 0 & 1 \\
ADNI3 & 36 & Siemens & Skyra & Accelerated & 55 & 2.0 & 0 & 0 & 1 & 0 & 1 \\
ADNI3 & 36 & Siemens & Skyra & Accelerated & 127 & 2.0 & 1 & 0 & 0 & 0 & 1 \\
ADNI3 & 94 & Siemens & Prisma & Accelerated & 55 & 2.0 & 1 & 0 & 0 & 0 & 1 \\
ADNI3 & 127 & GE & DISCOVERY MR750 & Accelerated & 127 & 2.0 & 0 & 0 & 1 & 0 & 1 \\
ADNI3 & 128 & GE & DISCOVERY MR750w & Accelerated & 36 & 2.0 & 1 & 0 & 0 & 0 & 1 \\
ADNI3 & 129 & GE & DISCOVERY MR750 & Accelerated & 22 & 2.0 & 0 & 0 & 0 & 1 & 1 \\
ADNI3 & 135 & GE & SIGNA Premier & Accelerated & 36 & 2.0 & 0 & 0 & 0 & 1 & 1 \\
ADNI3 & 168 & Siemens & Prisma\_fit & Accelerated & 42 & 2.0 & 0 & 0 & 1 & 0 & 1 \\
ADNI3 & 941 & Siemens & Prisma\_fit & Accelerated & 55 & 2.0 & 1 & 0 & 0 & 0 & 1 \\
ADNI4 & 389 & Siemens & Prisma\_fit & Accelerated & 127 & 2.0 & 0 & 0 & 0 & 1 & 1 \\
ADNI4 & 941 & Siemens & Prisma\_fit & Accelerated & 127 & 2.0 & 0 & 0 & 0 & 1 & 1 \\
ADNIGO/ADNI2 & 21 & GE & DISCOVERY MR750 & Accelerated & 46 & 2.7 & 0 & 1 & 0 & 0 & 1 \\
ADNIGO/ADNI2 & 27 & GE & DISCOVERY MR750w & Accelerated & 54 & 2.7 & 0 & 0 & 1 & 0 & 1 \\
ADNIGO/ADNI2 & 37 & Siemens & Prisma & Accelerated & 127 & 2.0 & 0 & 0 & 1 & 0 & 1 \\
ADNIGO/ADNI2 & 52 & GE & DISCOVERY MR750 & Accelerated & 46 & 2.7 & 0 & 0 & 0 & 1 & 1 \\
ADNIGO/ADNI2 & 52 & GE & DISCOVERY MR750 & Accelerated & 68 & 2.7 & 0 & 0 & 1 & 0 & 1 \\
ADNIGO/ADNI2 & 52 & GE & DISCOVERY MR750 & Accelerated & 127 & 2.7 & 0 & 0 & 1 & 0 & 1 \\
ADNIGO/ADNI2 & 57 & GE & DISCOVERY MR750 & Accelerated & 54 & 2.7 & 1 & 0 & 0 & 0 & 1 \\
ADNIGO/ADNI2 & 68 & Siemens & Verio & Accelerated & 55 & 3.0 & 1 & 0 & 0 & 0 & 1 \\
ADNIGO/ADNI2 & 98 & GE & DISCOVERY MR750 & Accelerated & 54 & 2.7 & 1 & 0 & 0 & 0 & 1 \\
ADNIGO/ADNI2 & 99 & GE & DISCOVERY MR750 & Unaccelerated & 54 & 5.0 & 1 & 0 & 0 & 0 & 1 \\
ADNIGO/ADNI2 & 109 & GE & Signa HDxt & Accelerated & 46 & 2.7 & 0 & 0 & 0 & 1 & 1 \\
ADNIGO/ADNI2 & 129 & GE & SIGNA HDx & Accelerated & 46 & 2.7 & 1 & 0 & 0 & 0 & 1 \\
ADNIGO/ADNI2 & 135 & GE & DISCOVERY MR750w & Accelerated & 36 & 2.7 & 0 & 0 & 1 & 0 & 1 \\
\end{longtable}
\footnotesize\textit{Note:} Each row is one observed combination of phase, site, scanner manufacturer and model, acceleration category, final NIfTI-derived diffusion-volume count, and slice thickness. The 84 observed cells comprised 0 cells containing all four groups, 12 containing three groups, 34 containing two groups, and 38 containing one group, demonstrating incomplete crossing of diagnosis and acquisition factors.
\end{landscape}
% End inlined acquisition_cells_R2.tex

\begin{landscape}
\begin{longtable}{p{0.25\linewidth}p{0.08\linewidth}p{0.43\linewidth}p{0.15\linewidth}}
\caption{Multiplicity families used in the revised analyses.}\label{tab:multiplicity_families}\\
\toprule Analysis family & Tests & Definition & Correction\\ \midrule
\endfirsthead
\toprule Analysis family & Tests & Definition & Correction\\ \midrule
\endhead
Aggregated networks & 16 & Eight networks, separately for average- and modal-rank endpoints & BH-FDR within endpoint (8)\\
Tian structure omnibus & 14 & Seven structures, separately for average- and modal-rank endpoints & BH-FDR within endpoint (7)\\
Primary hippocampal pairwise & 12 & Six contrasts for each primary rank endpoint & BH-FDR across 12\\
Raw-value omnibus & 2 & Raw average and modal controllability & BH-FDR across 2\\
Density-adjusted raw omnibus & 2 & Density-adjusted raw average and modal controllability & BH-FDR across 2\\
Raw-value pairwise & 12 & Six contrasts for each raw endpoint & BH-FDR across 12\\
Hippocampal long axis & 8 & Four subdivisions, separately by endpoint & BH-FDR within endpoint (4)\\
Cognition associations & 10 & Two rank-based controllability endpoints by five cognitive outcomes, adjusted for age, sex, and education & BH-FDR across 10\\
Matched FreeSurfer cognition & 15 & Three imaging measures by five cognitive outcomes & BH-FDR across 15\\
Topology group contrasts & 4 & Four exploratory contrasts for hippocampal betweenness & BH-FDR across 4\\
Topology--controllability & 4 & Two unranked endpoints by strength and betweenness & BH-FDR across 4\\
Incremental cognitive models & 5 & Five cognitive outcomes & BH-FDR across 5\\
Biomarker sensitivity & 16 & Two endpoints within biomarker, window, and tested-effect strata & BH-FDR within each 2-test family\\
Low-coverage exclusion sensitivity & 2 & Average- and modal-rank group effects after excluding five review-trigger cases & BH-FDR across 2\\
Low-dimensional classification & 24 & Six feature sets in four tasks & Holm across 24\\
Paired classification comparisons & 20 & Five candidate-versus-volume comparisons in four tasks & Holm across 20\\
\bottomrule
\end{longtable}
\footnotesize\noindent\textit{Note:} BH-FDR, Benjamini--Hochberg false discovery rate correction; Holm, Holm familywise error rate correction.
\end{landscape}

\begin{table}[!htbp]
\centering
\caption{Acquisition-overlap diagnostics and design-balanced sensitivity analyses.}
\label{tab:acquisition_overlap}
\footnotesize
\setlength{\tabcolsep}{3pt}
\textbf{A. Omnibus results and balance diagnostics}\par\vspace{1mm}
\begin{tabularx}{\textwidth}{@{}>{\raggedright\arraybackslash}p{0.23\textwidth}>{\raggedright\arraybackslash}Xccccc@{}}
\toprule
Analysis & Endpoint or diagnostic & $n$ or ESS & Statistic & $p$ & $p_{\mathrm{FDR}}$ & Partial $\eta^2$\\
\midrule
Strict covariate adjustment & Average rank & 281 & $F_{3,210}=4.368$ & 0.0057 & 0.0093 & 0.059\\
& Modal rank & 281 & $F_{3,210}=3.977$ & 0.0093 & 0.0093 & 0.054\\
\addlinespace
Generalized overlap weighting & Overall effective sample size & 45.2 & Max pairwise SMD: 1.640 to 1.042 & -- & -- & --\\
& Average rank & 45.2 & Wald $\chi^2_3=8.875$ & 0.0310 & 0.0310 & 0.098\\
& Modal rank & 45.2 & Wald $\chi^2_3=8.979$ & 0.0296 & 0.0310 & 0.099\\
\addlinespace
Connected exact-cell common support & Average rank & 218 & $F_{3,164}=5.663$ & 0.0014 & 0.0025 & 0.094\\
& Modal rank & 218 & $F_{3,164}=5.112$ & 0.0025 & 0.0025 & 0.086\\
\addlinespace
All-four-groups-per-cell restriction & Feasibility diagnostic & 0 & No exact cell contained all four groups & -- & -- & --\\
\bottomrule
\end{tabularx}
\par\vspace{3mm}
\textbf{B. Covariate-adjusted estimates in the connected exact-cell common-support sample}\par\vspace{1mm}
\begin{tabular}{llrrr}
\toprule
Endpoint & Group & $n$ & Adjusted mean & 95\% CI\\
\midrule
Average rank & CN & 93 & 345.73 & 325.44--366.01\\
& SMC & 24 & 383.14 & 340.41--425.87\\
& MCI & 66 & 321.88 & 297.12--346.63\\
& AD & 35 & 286.03 & 255.84--316.23\\
\addlinespace
Modal rank & CN & 93 & 107.18 & 86.21--128.16\\
& SMC & 24 & 69.41 & 25.86--112.95\\
& MCI & 66 & 131.06 & 105.22--156.91\\
& AD & 35 & 165.42 & 134.52--196.32\\
\bottomrule
\end{tabular}
\par\vspace{1mm}
\footnotesize{\textit{Note:} The strict model included 12 scanner-model, three phase, 43 site, two acceleration, nine NIfTI-derived diffusion-volume-count, and four slice-thickness levels, together with age, sex, education, and two motion measures; the nominal 76-column full design had rank 71 with 210 residual degrees of freedom. Group-specific overlap-weighting effective sample sizes were 34.8 CN, 11.5 SMC, 6.7 MCI, and 20.9 AD. The weighted partial-$\eta^2$ bootstrap 95\% CIs were 0.032--0.298 for average rank and 0.031--0.292 for modal rank. Residual balance remained inadequate, so weighted results are diagnostic irrespective of their nominal tests. The connected analysis retained 46 exact acquisition cells containing at least two groups and required direct shared-cell support for all six group pairs; permutations were restricted within cell. Direct SMC--MCI support comprised four cells with five SMC and six MCI participants. No exact cell contained all four groups. BH-FDR correction was applied across the two endpoints within each estimable analysis. SMD, standardized mean difference; ESS, effective sample size.}
\end{table}

\begin{table}[!htbp]
\centering
\caption{Diffusion-protocol distribution and retrospective registration/label-placement QC.}
\label{tab:protocol_qc_summary}
\footnotesize
\begin{tabularx}{\textwidth}{@{}p{0.34\textwidth}X@{}}
\toprule
Item & Result\\
\midrule
Software & MRtrix3 3.0.6; FSL 6.0.7.18, including FLIRT\\
NIfTI-derived total-volume categories & 22 (1), 31 (44), 36 (24), 42 (1), 46 (61), 54 (39), 55 (87), 68 (9), and 127 (15)\\
Observed b-value shell sets & $b=0/1000$, $b=0/1300$, and $b=0/500/1000/2000$~s/mm$^2$\\
Observed b0-count categories & 1, 4, 5, 6, 7, 8, and 13\\
Observed nonzero diffusion-weighted volume-count categories & 16, 30, 32, 36, 41, 48, 60, and 114\\
Protocol strata: $b=0/1000$ & 257 participants: b0/nonzero counts 6/16 (1), 1/30 (44), 4/32 (24), 6/36 (1), 5/41 (61), 6/48 (39), and 7/48 (87)\\
Protocol strata: $b=0/1300$ & 9 participants: b0/nonzero count 8/60\\
Protocol strata: $b=0/500/1000/2000$ & 15 participants: b0/nonzero count 13/114\\
Response/FOD commands & \texttt{dwi2response dhollander} and \texttt{dwi2fod msmt\_csd} were applied to every dataset\\
Connectome scaling & \texttt{-scale\_invnodevol}: $A_{ij}=2\sum_{s\in S_{ij}}w_s/(N_i+N_j)$, where $N_i$ and $N_j$ are node-label voxel counts\\
Phase-encoding metadata & BIDS directions $i$ (2), $j$ (244), and $j-$ (35) in the retrospective input audit; all 15 127-volume inputs were PA; reverse phase-encoding pairs were not used for susceptibility correction\\
Retrospective quantitative QC completed & 281/281 participants\\
Brain-mask Dice & Mean 0.917; 0/281 below 0.70\\
Hippocampal-label coverage in DWI brain mask & Mean 0.992; 5/281 below 0.90 (CN 1, SMC 0, MCI 2, AD 2)\\
Manual review of coverage-trigger cases & Acceptable T1--DWI registration and hippocampal-label placement confirmed in all five cases; no gross misregistration or erroneous hippocampal localization\\
Bilateral hippocampal labels present & 281/281 participants\\
QC montage generation & Registration, 281/281; hippocampal placement, 280/281; the one montage-generation failure was display-only and retained complete quantitative QC\\
Existing connectome availability and numerical checks & 281/281 matrices available; all passed dimension, finite-value, non-negativity, symmetry, and zero-diagonal checks\\
Additional exclusions during retrospective audit & Registration, 0; tractography/connectome availability, 0; QC thresholds, 0\\
Sensitivity excluding five coverage-trigger cases & $n=276$ (CN 109, SMC 27, MCI 86, AD 54); average rank $F_{3,269}=5.554$, $p_{\mathrm{FDR}}=0.0017$, partial $\eta^2=0.058$; modal rank $F_{3,269}=5.264$, $p_{\mathrm{FDR}}=0.0017$, partial $\eta^2=0.055$\\
\bottomrule
\end{tabularx}
\par\vspace{1mm}
\footnotesize{\textit{Note:} Counts in parentheses are participants in the final NIfTI input categories. All 127-volume inputs were PA acquisitions and each contained 13 b0 and 114 nonzero diffusion-weighted volumes. Dice and coverage are gross-overlap checks and do not prove anatomical localization accuracy. QC thresholds were review triggers rather than post hoc exclusion rules. The two exclusion-sensitivity tests formed one BH-FDR family. Candidate-level failure counts preceding formation of the final 281-participant analytic cohort were not available in the archived processing records.}
\end{table}
% End inlined supplementary_new_tables.tex

\FloatBarrier

\section{Supplementary Figures}

\begin{figure}[!htbp]
    \centering
    \includegraphics[width=0.9\linewidth]{SI_1.pdf}
    \caption{Descriptive cortical spatial distributions of average and modal controllability across clinically defined groups. (A) Average controllability and (B) modal controllability were averaged within each group and mapped to the Schaefer2018 cortical surface. A common color scale was used across CN, SMC, MCI, and AD within each metric to permit visual comparison. These maps are descriptive only; no parcel-level inference was performed, and visual differences should not be interpreted as statistically significant.}
    \label{fig:SI_figure1}
\end{figure}

\begin{figure}[!htbp]
    \centering
    \includegraphics[width=0.9\linewidth]{SI_2.pdf}
    \caption{Descriptive structural-connectivity patterns of the hippocampus and Tian atlas-defined non-cortical structures across clinically defined groups. (A) Major hippocampal connections retained for visualization. (B) Mean connectivity between the hippocampus and the visual (VN), somatomotor (SMN), dorsal attention (DAN), ventral attention (VAN), limbic (LN), frontoparietal (FPN), default mode (DMN), and atlas-defined non-cortical (NonCor) target groups. (C) Major connections between Tian atlas-defined structures and cortical networks. (D) Mean connectivity between Tian atlas-defined structures and cortical networks. A common display threshold was used only to improve visualization and was not applied to inferential analyses; between-group visual differences should therefore be interpreted descriptively.}
    \label{fig:SI_figure2}
\end{figure}

\begin{figure}[!htbp]
    \centering
    \setlength{\tabcolsep}{2pt}
    \begin{tabular}{cccc}
    \textbf{CN} & \textbf{SMC} & \textbf{MCI} & \textbf{AD}\\
    \includegraphics[width=0.235\linewidth]{qc_overlays/qc_CN_registration.png} &
    \includegraphics[width=0.235\linewidth]{qc_overlays/qc_SMC_registration.png} &
    \includegraphics[width=0.235\linewidth]{qc_overlays/qc_MCI_registration.png} &
    \includegraphics[width=0.235\linewidth]{qc_overlays/qc_AD_registration.png}\\
    \includegraphics[width=0.235\linewidth]{qc_overlays/qc_CN_hippocampus.png} &
    \includegraphics[width=0.235\linewidth]{qc_overlays/qc_SMC_hippocampus.png} &
    \includegraphics[width=0.235\linewidth]{qc_overlays/qc_MCI_hippocampus.png} &
    \includegraphics[width=0.235\linewidth]{qc_overlays/qc_AD_hippocampus.png}
    \end{tabular}
    \caption{Representative retrospective registration and hippocampal-label-placement QC overlays. Columns show one CN, SMC, MCI, and AD participant selected near the cohort median brain-mask Dice value with complete hippocampal coverage. The top row shows the transformed T1-derived brain-mask contour over the mean b0 image; the bottom row shows transformed bilateral hippocampal-label contours over the mean b0 image. These examples illustrate the QC display, whereas quantitative registration and label-coverage metrics were computed for all 281 participants.}
    \label{fig:registration_qc_overlays}
\end{figure}

\begin{figure}[!htbp]
    \centering
    \includegraphics[width=0.82\linewidth]{qc_overlays/hippocampus_QC.png}
    \caption{Example hippocampal-label overlay from a review-trigger case in the CN group. The transformed bilateral hippocampal-label coverage within the DWI brain mask was 0.894. Manual inspection confirmed acceptable T1--DWI registration and hippocampal-label placement without gross misregistration or erroneous hippocampal localization. The coverage threshold was used to trigger review rather than automatic exclusion; the primary group effects were retained after excluding all five coverage-trigger cases (Supplementary Table~S21).}
    \label{fig:low_coverage_qc_example}
\end{figure}

\FloatBarrier
\end{document}
